\section{The single-coin obstruction survives uniformly small bias}
\label{sec:small-bias-obstruction}
\noindent\textbf{PROVED-ON-PAPER; INDEPENDENT VERIFICATION/REFINEMENT.}
The small-bias extension was already present in the concurrently maintained
critical-two research notes when the separate reviewer first inspected them.
This independent derivation is not counted as discovery of a new theorem.
It is a single-coin obstruction, not a theorem about unrestricted population
algorithms. The finite-gap sign entropy lemma below is the sole imported
mathematical input. Its separate audit is part of the companion review.
This section is not Lean formalized or externally peer reviewed.

Fix $a=1/4$, $\mu(dd)=2d/a^2\,dd$ on $(0,a]$,
$0<\rho\le1/4$, $L=1+\log(1/\rho)$ and $H=1+\log L$.
One genuine coin has hidden $D\sim\mu$ and
$S\mid D=d\sim\operatorname{Bernoulli}(q(d))$; its observations are
conditionally iid $\operatorname{Bernoulli}(1/2+(2S-1)D)$.
Let $\mathcal Q_\rho=\{q:0\le q\le1,\int q\,d\mu\le\rho\}$ and
$\theta(q)=\int q\,d\mu$. An observable single-coin procedure may use an
independent seed, output any real $Y$, and stop at $T=0$ or at an unbounded
stopping time. Only its genuine coin flips are charged. The seed and policy
cannot depend on $q$ or the realized hidden gap.

\begin{theorem}\label{sb:theorem}
There is a universal $c>0$ such that every such procedure with
\[
 V:=\sup_{q\in\mathcal Q_\rho}\operatorname{Var}_q(Y)<\infty,
 \qquad W:=\sup_{q\in\mathcal Q_\rho}\mathbb E_qT<\infty,
 \qquad
 \sup_{q\in\mathcal Q_\rho}|\mathbb E_qY-\theta(q)|\le\epsilon
 \le\rho/64
\]
satisfies $VW\ge c\rho LH$. The same lower bound holds for
$MW$, where $M:=\sup_q\mathbb E_qY^2$.
\end{theorem}

\paragraph{The finite-gap input.}
The following statement is used, with the common randomized test independent
of the unknown gap. If
$0<d_K<\cdots<d_1\le1/128$, $d_j/d_{j+1}\ge2$,
$K\ge\lceil e^{32}\rceil$, and one observable sign test has error at most
$1/10$ at every $1/2\pm d_j$ and finite expected stopping time at every one
of these $2K$ laws, then
\begin{equation}\label{sb:entropy}
 \sum_{j=1}^K d_j^2\mathbb E_{-d_j}\tau
 \ge\frac1{512}K\log K.
\end{equation}
No stopping or integrability is required at the fair-coin law. The companion
single-coin proof establishes this from finite-prefix KL, time-window
overlap, and the binary data-processing inequality. In particular this is
not a direct appeal to a pointwise or limsup version of a sign-testing bound.

\begin{proof}[Proof of Theorem~\ref{sb:theorem}]
First $M\le V+(\rho+\epsilon)^2<\infty$. Write
$m_s(d)=\mathbb E_{d,s}Y$ and $M_s(d)=\mathbb E_{d,s}Y^2$ on a common
full-measure set. Such a set exists because the laws $q=0$ and $q=\rho$
make both conditional second moments integrable. Set $g=m_+-m_-$ and
$b_q=\mathbb E_qY$. For every feasible $q$,
\begin{equation}\label{sb:response}
 \left|\int q(g-1)\,d\mu\right|
 =|b_q-b_0-\theta(q)|\le2\epsilon.
\end{equation}
Taking $q=\rho$ gives $\int g\,d\mu\ge31/32$. Conditional variance under
this constant sign allocation and Cauchy--Schwarz imply
\begin{equation}\label{sb:variance}
 V\ge\rho(1-\rho)\int g^2\,d\mu
 \ge\frac{2883}{4096}\rho>\frac7{10}\rho.
\end{equation}
Indeed conditioning on $D$ and then the two signs contributes
$\rho(1-\rho)g(D)^2$ to the variance. Since
$(\rho+\epsilon)^2\le(4225/16384)\rho$, it follows that
\begin{equation}\label{sb:M-V}
 M\le\tfrac32 V.
\end{equation}

We next remove artificial thinning before the first observation. The event
$A=\{T\ge1\}$ is determined by the seed alone, so its probability $p$ is
the same for every $q$ and every fixed signed gap. If $p=0$, the expectation
would be constant in $q$, contradicting the bias guarantee at
$\theta=0$ and $\theta=\rho$; hence $p>0$. Let
$c_0=\mathbb E[Y\mathbf1_{A^c}]$, also independent of $q$.
Condition the seed on $A$, run the original procedure, and output
$Z=pY+c_0$. Resampling seeds until $A$ occurs implements this conditioning
almost surely using free independent randomness; the constants $p,c_0$
are fixed by the original policy and may be granted in this lower-bound
argument. The new procedure always takes at least one genuine flip,
preserves each expectation $b_q$, and has
$W_Z=W/p\ge1$.
For $p<1$, Cauchy--Schwarz on $A^c$ gives
\begin{align*}
 \mathbb E_qZ^2
 &=p\mathbb E_q[Y^2\mathbf1_A]+2c_0b_q-c_0^2\\
 &\le pM+2c_0b_q-\frac{c_0^2}{1-p}
 \le pM+(1-p)b_q^2
 \le pM+(\rho+\epsilon)^2.
\end{align*}
For $p=1$ the last bound is immediate. Apply \eqref{sb:variance} also to
$Z$, whose bias bound is unchanged. The last display implies
$pM\ge(7/10-4225/16384)\rho\ge2\rho/5$.
Set
\[
 P=MW,\qquad
 m=\max\{\rho,pM+(\rho+\epsilon)^2\}.
\]
The preceding inequalities imply
\begin{equation}\label{sb:normalization}
 \rho\le m\le\tfrac52pM\le\tfrac52P,
 \qquad mW_Z\le\tfrac52P,
 \qquad P\ge\tfrac25\rho.
\end{equation}
In what follows $m_s,M_s,w_s$ denote the conditional means, second moments,
and expected work of the active procedure $Z$. They are finite for both
signs on a common full-measure set. Its uniform second moment is at most
$m$, and \eqref{sb:response} remains valid.

\paragraph{Good gaps without exact unbiasedness.}
For $E=\{g<7/8\}$, nonatomicity allows a subset of mass exactly $\rho$
if $\mu(E)\ge\rho$. Its indicator would contradict
\eqref{sb:response}, since $\rho/8>2\epsilon$.
Thus $\mu(E)<\rho$; applying \eqref{sb:response} to $q=\mathbf1_E$
now gives
\[
 \mu(E)\le16\epsilon\le\rho/4.
\]
The other bad sets are
$E_-:=\{M_->1/64\}$ and
$E_+:=\{M_+-M_->2m/\rho\}$.
Markov's inequality under $q=0$ gives $\mu(E_-)\le64m$.
If $\mu(E_+)\ge\rho$, a subset $F$ of mass $\rho$ and
$q=\mathbf1_F$ would give
$\mathbb E_qZ^2=\int M_-\,d\mu+\int_F(M_+-M_-)\,d\mu>2m$,
contradicting the uniform second-moment bound. Hence $\mu(E_+)<\rho$.
On the complement $G$ of these sets and the common null set,
\begin{equation}\label{sb:good}
 \mu(G^c)\le66m,\quad
 |m_-|\le1/8,\quad m_+\ge3/4,\quad
 M_-\le1/64,\quad M_+\le1/64+2m/\rho.
\end{equation}

Run the active procedure
$n=\lceil256(1+m/\rho)\rceil\le513m/\rho$ times, with independent
seeds on successive unused blocks of the same fixed-gap coin. Compare the
average output to $7/16$. At every gap in $G$ the distance of either mean
from this threshold is at least $5/16$. Conditional independence of
successive blocks follows by summing over the finite stopping prefixes;
their conditional finite expected lengths imply almost-sure termination.
Chebyshev's inequality gives, for either sign, an error at most
\[
 \frac{256}{25n}\left(\frac1{64}+\frac{2m}{\rho}\right)
 \le\frac2{25}<\frac1{10}.
\]
The negative expected work of this common test is $nw_-(d)$.

\paragraph{Select separated gaps and apply the entropy bound.}
For $j\ge3$, put
$I_j=(a4^{-j}/2,a4^{-j}]$, so
$\mu(I_j)=(3/4)16^{-j}$. Set
\[
 J=\left\lfloor\log_{16}\frac1{176m}\right\rfloor,
 \qquad K=J-2.
\]
When $K\ge\lceil e^{32}\rceil$, every $j=3,\ldots,J$ satisfies
$\mu(I_j)\ge132m$, so at least half its mass is in $G$.
Choose $d_j\in G\cap I_j$ with
\[
 w_-(d_j)\le\frac4{\mu(I_j)}\int_{I_j}w_-\,d\mu.
\]
Such points exist by averaging; conditional finite moments and work may
also be imposed because the discarded set is null. These gaps have ratios
strictly greater than two and maximum at most $1/256$. Moreover
\[
 \sum_{j=3}^J d_j^2w_-(d_j)
 \le\frac{16a^2}{3}\sum_{j=3}^J\int_{I_j}w_-\,d\mu
 \le\frac{16a^2}{3}W_Z.
\]
Applying \eqref{sb:entropy} and \eqref{sb:normalization},
\begin{equation}\label{sb:entropy-product}
 \frac1{512}K\log K
 \le\frac{16a^2}{3}nW_Z
 \le6840a^2\frac P\rho.
\end{equation}

It remains to compare scales without assuming the variance is small.
If $P\ge\sqrt\rho$, then $P\ge c_1\rho LH$ for a universal
$c_1>0$, since $\sqrt\rho LH$ is bounded on $(0,1/4]$.
If $P<\sqrt\rho$, then $m<(5/2)\sqrt\rho$, and therefore
\[
 K\ge\frac{\tfrac12\log(1/\rho)-\log440}{\log16}-3.
\]
For all sufficiently small universal $\rho_0>0$, this ensures the
cardinality required in \eqref{sb:entropy}, $K\ge c_2L$, and
$\log K\ge c_3H$. Equation~\eqref{sb:entropy-product} then gives
$P\ge c_4\rho LH$. For $\rho\in[\rho_0,1/4]$ use
$P\ge2\rho/5$ and boundedness of $LH$ on that fixed compact interval.
This proves the universal lower bound for $MW$. Finally
\eqref{sb:M-V} yields the asserted bound for $VW$.
\end{proof}

\begin{corollary}\label{sb:arithmetic-average}
Suppose an estimator is the deterministic arithmetic average of $n\ge1$
independent copies of the same single-coin procedure, and its uniform
expected work is at most $B$. If its worst-case RMS error is at most
$\eta\le\rho/64$, then $B\eta^2\ge c\rho LH$.
The procedure and $n$ may depend on $B,\rho$.
\end{corollary}
\begin{proof}
For every $q$ the squared error of the average is
$\operatorname{Var}_q(Y)/n+(\mathbb E_qY-\theta(q))^2$.
Thus its kernel has bias at most $\eta$ and $V/n\le\eta^2$.
Uniform expected work gives $nW\le B$, so
$B\eta^2\ge(nW)(V/n)=VW\ge c\rho LH$.
\end{proof}

The last corollary does not cover clipping or any other nonlinear
postprocessing, unequal kernels whose biases cancel, data-dependent
replication counts, or interaction among genuine reservoir coins.
In particular, a raw variance lower bound cannot be substituted for a
lower bound on clipped probability-estimation risk. The unrestricted
population minimax problem remains open in this pass.
